\section{Introducción: Análisis de la Incidencia Delictiva Estatal}

No es noticia nueva de que México sufre de un serio problema de delicuencia. Y esto, desgraciadamente,
no ha disminuido considerablemente en 10 años a la fecha. Pero la realidad es que distntos estados
de la república sufren de delitos en particular que, aunque presentes en otros lugares, no son el
principal foco del lugar.

Nuestro objetivo al analizar la incidencia delectiva es poder identificar no solo cuales son éstos
delitos los más preliminares en cada estado, si no también poder crear un modelo que establezca cual
es delito más probable que pueda sufrir un ciudadano de un estado en particular.

Desgraciadamente, todxs somos propensos a sufrirlos, pero identificarlos no solo nos permite a nostros
como habitantes poder prepararnos para evitarlos lo más que se pueda, como le puede servir a las autoridades
componentes para concientizar y, utópicamente, idear soluciones que reduzcan verdaderamente la epidemia
de delicuencia que sufre el país.

Este dataset pertenece a la \href{https://www.datos.gob.mx/}{Plataforma Nacional de Datos Abiertos Mexicana} y fue
creado por el Secretariado Ejecutivo del Sistema Nacional de Seguridad Pública (SESNSP). Este dataset, nombrado
simplemente como \textit{Incidencia Delictiva}, puede descargarse siguiendo \href{https://www.datos.gob.mx/dataset/incidencia_delictiva/resource/d9b2792a-33a2-4ea8-8527-210d9e99de5e}{este enlace}.
